% Isolated exact-reduction candidate. No nonvanishing or rank theorem asserted.
\subsection{The residual-pencil boundary underlying the surface cup product}

\begin{lemma}[A symmetric-square resolution and its actual cup boundary]
\label{lem:DE-A3-residual-pencil-boundary-candidate}
\status{Candidate; complete proof submitted for independent audit.}
Keep the sharp-source hypotheses and notation of the proved monad and
restricted-Petri lemmas. Put $\mathcal L=LP$, $E=Q^\vee$,
$G=\mathfrak g^\vee$, and $R=R_s=\omega_D A^{-1}$.
The dual of $0\to\mathfrak g\to Q\to i_*A\to0$ is
\[
0\longrightarrow E\xrightarrow{\varphi}G
\longrightarrow i_*R\longrightarrow0.
\tag{DE-RB1}
\]
Its boundary identifies $H^0(D,R)$ with $U=H^1(T,E)$.
There is an exact complex
\[
0\longrightarrow\mathcal L^{-1}
\xrightarrow{d_2}E\otimes G\xrightarrow{d_1}
\operatorname{Sym}^2G\longrightarrow i_*(R^2)
\longrightarrow0,
\tag{DE-RB2}
\]
where $d_1(e\otimes g)=\varphi(e)g$ and
$d_2(e\wedge e')=e\otimes\varphi(e')-e'\otimes\varphi(e)$.
Let $\partial_s:H^0(D,R^2)\to H^2(T,\mathcal L^{-1})$
be its two-step cohomological boundary. Under (DE-RB1), the
actual surface cup-wedge map is
\[
\mu=\partial_s\circ\operatorname{mult}:
\operatorname{Sym}^2H^0(D,R)\longrightarrow H^0(D,R^2)
\longrightarrow H^2(T,\mathcal L^{-1}).
\tag{DE-RB3}
\]
This holds also when the Cartier divisor $D$ is nonreduced.
\end{lemma}
\begin{proof}
The earlier local matrix for $f:\mathfrak g\to Q$ is
$\operatorname{diag}(uw,1)$ with $u$ a unit and $w$ a local
equation for $D$. Its dual has the same Smith form. Duality for
the Cartier divisor gives
$\mathcal E xt^1_T(i_*A,\mathcal O_T)=i_*(A^{-1}\mathcal L|_D)
=i_*R$, proving (DE-RB1). The acyclicity of $G$ makes its
cohomology boundary the stated isomorphism.

Exactness of (DE-RB2) is local. After absorbing the unit, write
$\varphi(e_1)=wf_1$, $\varphi(e_2)=f_2$. The four images under
$d_1$ are
\[
wf_1^2,\quad wf_1f_2,\quad f_1f_2,\quad f_2^2.
\]
Since $w$ is a nonzerodivisor, their relation module is freely
generated by $e_1\otimes f_2-w e_2\otimes f_1$, the image
of $d_2$. The cokernel is $(\mathcal O_T/(w))f_1^2$, precisely
the square of the quotient line $R$. This verifies the complex
over every local ring, without assuming $D$ reduced. The first
term is $\det E=\mathcal L^{-1}$.

We check the boundary, including diagonal products in
characteristic two. On an affine cover choose local lifts
$u_i,v_i$ in $G$ of $r,s\in H^0(D,R)$. Write
$u_j-u_i=\varphi(x_{ij})$ and
$v_j-v_i=\varphi(y_{ij})$. Then $x,y$ represent their
boundaries in $H^1(E)$. The products $u_iv_i$ lift $rs$ to
$\operatorname{Sym}^2G$. Their Cech difference lifts through
$d_1$ as
\[
z_{ij}=x_{ij}\otimes v_j+y_{ij}\otimes u_i.
\]
The next Cech difference is
\[
z_{jk}-z_{ik}+z_{ij}
=-x_{ij}\otimes\varphi(y_{jk})
  +y_{jk}\otimes\varphi(x_{ij})
=-d_2(x_{ij}\wedge y_{jk}).
\]
In characteristic two the sign is positive. The resulting
class is exactly the cup product followed by the alternating
bundle pairing. This proves (DE-RB3), for $r=s$ as well as
for distinct sections. In particular no symmetric/divided-power
identification or division by two has been used.
\end{proof}

\begin{corollary}[The precise residual-pencil restriction]
\label{cor:DE-A3-residual-pencil-restriction-candidate}
\status{Candidate equivalent reformulation; no nonvanishing asserted.}
At a sharp source, $H^0(D,R)$ has dimension two. If $D$ is
integral, its symmetric-square multiplication into $H^0(D,R^2)$
is injective and has dimension-three image $W_s$.
The variation family has image dimension
\[
\dim\operatorname{im}\Psi_s
=\dim\operatorname{im}
\left[W_s\xrightarrow{\partial_s}H^0(T,\mathcal L)^*
\longrightarrow\Delta_s^*\right].
\tag{DE-RB4}
\]
Thus the missing lower bound concerns this two-step boundary
and its actual determinant-direction restriction, even when
the multiplication itself is injective.
\end{corollary}
\begin{proof}
Take independent sections $r_0,r_1$ of $R$. On an integral
curve their ratio is nonconstant. A linear relation among
$r_0^2,r_0r_1,r_1^2$ would give a nonzero quadratic polynomial
over the algebraically closed base field satisfied by that
ratio, making it constant. This proves injectivity.
Equation (DE-RB4) is (DE-RB3) followed by the already proved
restricted Serre-duality identity (DE-PT4).
\end{proof}

\begin{remark}[The unresolved boundary kernel]
\label{rem:DE-A3-residual-pencil-boundary-scope-candidate}
\status{Candidate scope statement; the primitive positive-index gate remains open.}
The new identification involves the residual pencil $R_s$ and
the explicit boundary of (DE-RB2). It is not the ordinary
$H^0(A)\otimes H^0(R_s)$ multiplication map. Neither the
acyclicity of $G$ nor injectivity of multiplication on an
integral $D$ controls the kernel of $\partial_s$ after
restriction to $\Delta_s$. No variation-image lower bound,
common-radical vanishing, or new generic-rank theorem is claimed.
\end{remark}
